A space shuttle (N) orbits the Earth in the equatorial plane on a circular
trajectory with radius $r$. From the spaceship. The shuttle has an arm
designated to place satellites on the orbit. The arm is a metallic rod
(negligible mass) with length $l \ll r$. The arm is connected to the shuttle
with frictionless mobile articulation. A satellite \textbf{S} is attached to
the arm and let out from the shuttle. At a certain moment the angle between the
rod and the shuttle's orbit radius is $\alpha$, see the figure. You know the
mass of the Earth -- $M$, and the gravitational constant $G$.

\ioaafig{2014_TH_L2-f1.pdf}

\begin{description}
\item[(a)] Find out the values of angle $\alpha$ for which the configuration
  of the system shuttle--rod--satellite remains unchanged regardless to Earth
  (the system is in equilibrium), during orbiting the Earth. For each found
  value of angle $\alpha$, specify the type of the system equilibrium i.e.
  stable or unstable.

  You will take in to account the following assumptions: the initial orbit of
  the shuttle is not affected by the presence of the satellite S, all the
  external friction-type interactions are negligible, the satellite--shuttle
  gravitational interaction is negligible too. The following data are known:
  $m_1$ -- the mass of the shuttle, the mass of the satellite $m_2 \ll m_1$.
\item[(b)] In the moment of one stable equilibrium configuration, the rod with
  the satellite attached is slightly rotated with a very small angle in the
  orbital plane and then released. Demonstrate that the small oscillations of
  the satellite S relative to the shuttle are harmonically ones. Express the
  period $T_0$ of this cosmic pendulum as a function of the orbiting period $T$
  of the shuttle around the Earth.

  It is known the linear harmonic oscillator equation:
  \[ \frac{\mathrm{d}^2 \beta}{\mathrm{d}t^2} + \omega_0^2 \beta = 0; \quad
     \omega_0 = \frac{2\pi}{T_0}, \]
  where $\beta$ -- the instantaneous angular deviation, $T_0$ -- the period of
  the linear harmonic oscillator.
\item[(c)] If we consider that the mass of the satellite S, $m_2$, is not
  negligible by comparison with the shuttle's one $m_1$ in the conditions from
  point (a) the evolution on the orbit of the shuttle would be influenced by
  the presence of the satellite S rigid attached to the shuttle by the rod.
  Identify and determine the consequences on the shuttle's movement after one
  complete rotation around the Earth.
\item[(d)] Propose a special technical maneuver which can cancels the
  influence of the non negligible mass satellite S on the shuttles movement.
\end{description}
